\chapter{Protocols II: Binary Exchange Protocols}\label{ch:ll:protocols-ii-binary-exchange-protocols}

A futures exchange sends every packet of its market data twice, over two separate networks, as binary messages at fixed
offsets that a program reads without parsing. CME Group's documentation tells its clients to ``process both the Incremental
Feed A and Incremental Feed B due to the unreliable nature of UDP transport'', to keep each sequence number from whichever
line brings it first, and to treat a sequence number that neither brought as a loss on both lines, which calls for
recovery. The feed handler's first job is therefore not to understand the market but to keep the first copy of everything
and forget the second. This chapter reads the binary protocols of the exchange simulator of One Quant Book~10, which follow
the public Nasdaq specifications: order-by-order messages at fixed offsets, packets that carry sequence numbers over
\omterm{def:ll:protocols-ii-binary-exchange-protocols:multicast}{multicast}, snapshot and retransmission channels for recovery, and order entry over a sequenced session. It generates the
codecs from a schema instead of writing them by hand, arbitrates two lossy lines, and measures what decoding costs.

\section{Order-by-order feeds and their messages}

A venue's direct feed (One Quant Book~1, chapter~9) publishes every event of its book: an order added, executed, reduced,
deleted, replaced, a trade against hidden liquidity, an auction's cross. This is the market-by-order view of One Quant
Book~1, chapter~19, from which a subscriber rebuilds the book to any depth. Nasdaq's TotalView-ITCH is the model the
simulator follows: each message starts with a one-byte type, then fields at fixed offsets, with ``all integer fields \dots\
big endian (network byte order)''. An Add Order is 36 bytes: the type at offset 0, a two-byte stock locate code at 1 (a small
integer ``employed with the intent of serving as an array index''), a tracking number at 3, a six-byte timestamp in
nanoseconds since midnight at 5, the order reference at 11, the side at 19, the shares at 20, the symbol at 24 and the price
at 32, an integer in ten-thousandths (\cref{fig:ll:protocols-ii-binary-exchange-protocols:layout}). Nothing needs to be
searched: the price of an Add Order is always the four bytes at offset 32.

\begin{omfigure}
\begin{tikzpicture}[font=\scriptsize,b/.style={draw=gray!70,minimum height=0.55cm,inner sep=1pt,anchor=west,align=center}]
\def\u{0.33}
\node[anchor=west,font=\footnotesize\bfseries] at (0,1.0) {MoldUDP64 packet};
\node[b,fill=gray!12,minimum width=10*\u cm] (s) at (0,0.35) {session (10)};
\node[b,fill=omDef!25,minimum width=8*\u cm] (q) at (s.east) {sequence number (8)};
\node[b,fill=omDef!12,minimum width=2*\u cm] (c) at (q.east) {count};
\node[b,fill=gray!12,minimum width=2*\u cm] (l) at (c.east) {len};
\node[b,fill=omThm!15,minimum width=7*\u cm] (m) at (l.east) {message 1};
\node[b,fill=gray!12,minimum width=2*\u cm] (l2) at (m.east) {len};
\node[anchor=west,font=\footnotesize\bfseries] at (0,-0.55) {Add Order message (36 bytes)};
\foreach \w/\c/\t [count=\k] in {1/omThm!35/A,2/gray!12/loc,2/gray!12/trk,6/omMeth!20/time (ns),8/omDef!20/reference,
  1/omThm!35/B,4/omMeth!20/shares,8/gray!12/stock,4/omDef!20/price} {
  \ifnum\k=1 \node[b,fill=\c,minimum width=\w*\u cm,font=\tiny] (f\k) at (0,-1.2) {\t};
  \else \pgfmathtruncatemacro{\j}{\k-1}\node[b,fill=\c,minimum width=\w*\u cm,font=\tiny] (f\k) at (f\j.east) {\t}; \fi
}
\foreach \k/\o in {1/0,2/1,3/3,4/5,5/11,6/19,7/20,8/24,9/32} {
  \node[font=\tiny,anchor=north] at ([yshift=-0.02cm]f\k.south west) {\o};
}
\node[font=\tiny,anchor=north] at ([yshift=-0.02cm]f9.south east) {36};
\node[anchor=west,font=\scriptsize] at (0,-2.2) {byte offsets below each field; integers big-endian; price in 1/10\,000};
\end{tikzpicture}

\omcaption{A MoldUDP64 packet (session, the sequence number of its first message, a message count, then length-prefixed
message blocks) and the 36-byte Add Order message it carries, with the byte offset of each field. The layout is the one of
Nasdaq's specifications, which the simulator of One Quant Book~10 follows.}\label{fig:ll:protocols-ii-binary-exchange-protocols:layout}
\end{omfigure}

\begin{definition}[Flyweight codec]\label{def:ll:protocols-ii-binary-exchange-protocols:flyweight}
A \emph{flyweight codec}\index{flyweight codec} reads and writes a message in place: a small object that holds only a
pointer to the message's bytes and decodes a field from its fixed offset when the field is asked for, so that decoding a
message copies nothing and costs nothing for the fields not read.
\end{definition}

The build generates one flyweight per message from Book~10's machine-readable schema, a reader and a writer, and a dispatcher
that calls a visitor with the right type for a message's first byte (\cref{lst:ll:protocols-ii-binary-exchange-protocols:flyweight}).
Book~1's decoder of chapter~28 did the opposite: it copied every message into a common structure and called a
\texttt{std::function} for it. Both are correct; \cref{fig:ll:protocols-ii-binary-exchange-protocols:decode} measures the
difference.

\omcode{code/firm/wirecodec/cpp/firm_wirecodec.hpp}{39}{51}{The generated flyweight of the Add Order message: a pointer, and
a byte swap at a fixed offset for each field read.}\label{lst:ll:protocols-ii-binary-exchange-protocols:flyweight}

\section{Packets, sequence numbers and multicast}

\begin{definition}[Multicast, message gap]\label{def:ll:protocols-ii-binary-exchange-protocols:multicast}
\emph{Multicast}\index{multicast} is IP delivery of one packet to every host that has joined a group address: the sender
transmits once, and the network copies the packet towards each subscriber. A \emph{message gap}\index{message gap} is a range
of sequence numbers that a receiver has not received when a later one arrives.
\end{definition}

Market data is published over UDP \omterm{def:ll:protocols-ii-binary-exchange-protocols:multicast}{multicast} because one packet then reaches every subscriber at the same time, whatever their
number: MoldUDP64, the packet protocol of the simulator's feed, is described by its specification as a ``lightweight
protocol layer built on top of UDP'' in which ``each outbound packet is transmitted only once regardless of the number of
listeners''. UDP delivers without acknowledgement, retransmission or ordering, so the protocol above it carries what the
receiver needs to notice a loss: each packet's header gives the sequence number of its first message and the number of
messages it carries (\cref{fig:ll:protocols-ii-binary-exchange-protocols:layout}). A packet whose first number is above the
next expected one reveals a gap of the difference. When there is no data, the publisher sends \omterm{def:ll:protocols-i-fix:session}{heartbeats}, packets with a count
of zero carrying the next sequence number, ``typically \dots\ once per second'', so that a loss is noticed even in a quiet
market; a count of 65\,535 marks the end of the session. The sequence numbers are those of One Quant Book~1, chapter~28, on a
wire.

\section{Incremental and snapshot channels, retransmission and recovery}

\begin{definition}[Incremental channel, snapshot channel]\label{def:ll:protocols-ii-binary-exchange-protocols:channels}
An \emph{incremental channel}\index{incremental channel} carries the feed's messages, each a change to the book, in sequence.
A \emph{snapshot channel}\index{snapshot channel} publishes, at a regular interval, the full state of each instrument's book
together with the incremental sequence number it reflects, so that a receiver that has lost messages, or just started, can
rebuild the book without replaying the day.
\end{definition}

A gap leaves the receiver two ways back. Retransmission asks a server for the missing sequence numbers: MoldUDP64 has a
request packet for this, the simulator answers up to 1\,000 messages per request from a window of its last 100\,000, and a
small gap is closed in a round trip. A gap beyond the window, or a receiver starting in the middle of the day, needs the
\omterm{def:ll:protocols-ii-binary-exchange-protocols:channels}{snapshot channel}: the simulator publishes, for each instrument, a begin marker with the incremental sequence number the
snapshot reflects, the trading state, every displayed order in priority order, and an end marker with a checksum. The
receiver keeps buffering incremental messages, waits for the next snapshot, loads it, discards the buffered messages the
snapshot already reflects and applies the rest (the snapshot-and-delta pattern of One Quant Book~3, chapter~26, and the
snapshot recovery of One Quant Book~10, chapter~26). CME's snapshot messages carry the same synchronisation field,
\emph{LastMsgSeqNumProcessed}, ``used to synchronize the snapshot loop with the real-time feed''. The recovery time is set
by the snapshot interval, and \cref{prop:ll:protocols-ii-binary-exchange-protocols:both} says how often it is needed.

\section{Arbitration of redundant lines}

\begin{definition}[Line arbitration]\label{def:ll:protocols-ii-binary-exchange-protocols:arbitration}
\emph{Line arbitration}\index{line arbitration} merges the identical packets of redundant feed lines (One Quant Book~10,
chapter~26) into one stream: each sequence number is processed from whichever line delivers it first and discarded when the
other line delivers it again, and a gap is declared only when a later sequence number arrives before any line has delivered
the missing ones.
\end{definition}

\omcode{code/low-latency/16-protocols-ii-binary-exchange-protocols/cpp/ll_arb.hpp}{39}{57}{The arbiter: a duplicate is
dropped, a jump forward is a gap lost on both lines, everything else is new.}\label{lst:ll:protocols-ii-binary-exchange-protocols:arb}

Redundancy multiplies probabilities only when the lines fail independently, which is the point of sending them over separate
networks, and the reason the simulator draws each line's losses from its own random stream.

\begin{proposition}[Losing a packet on both lines]\label{prop:ll:protocols-ii-binary-exchange-protocols:both}
Let each line lose a given packet with probability $p$, independently of the other, and let a common cause (a shared switch,
the publisher itself) lose it on both with probability $c$, independently of the lines' own losses. The arbitrated stream
loses the packet with probability $c + (1 - c)p^2$. If each line's losses follow a two-state Gilbert--Elliott process that
moves from the good to the bad state with probability $g$ per packet, back with probability $h$, and loses a packet with
probability $\ell$ in the bad state and none in the good one, the stationary loss rate of a line is $p = \ell\, g/(g+h)$.
\end{proposition}

\begin{proof}
The packet is lost on the arbitrated stream if the common cause strikes, or if it does not and both lines lose it; the latter
events are independent, with probability $p^2$. For the two-state chain, the stationary probability $\pi$ of the bad state
balances the flows between states, $\pi h = (1-\pi) g$, so $\pi = g/(g+h)$, and a packet is lost with probability $\ell\pi$.
\end{proof}

Burstiness does not change a line's average loss rate, but it does change what a receiver sees: losses arrive in runs,
several packets at once, too many for a quick retransmission at the worst moment. And a common cause, however rare, dominates:
at $p = 10^{-4}$, independent losses on both lines cost $10^{-8}$ per packet, and a shared failure of $10^{-6}$ is a hundred
times more.

\Cref{fig:ll:protocols-ii-binary-exchange-protocols:gaps} runs the arithmetic on the simulator. Sixty seconds of a busy
instrument, about 32\,000 packets, with each line losing 0.1\% of packets independently plus Gilbert--Elliott bursts, and a
60-millisecond outage on each line, the two overlapping by 25 milliseconds. Line A shows 56 gaps and 160 missing messages, line
B 65 gaps and 186; the arbitrated stream shows two: the 17 messages published while both lines were out, and one packet lost on
line B while line A was down.

\begin{omfigure}
\begin{tikzpicture}
\begin{axis}[width=12cm,height=4.8cm,xmin=0,xmax=61,ymin=0.4,ymax=3.6,ytick={1,2,3},
  yticklabels={arbitrated,line B,line A},xlabel={seconds after the open},tick label style={font=\small},
  label style={font=\small},grid=major,grid style={gray!25},
  scatter/classes={1={omThm},2={omDef},3={omDef}},table/col sep=comma]
\addplot[only marks,mark=|,mark size=4pt,omDef,thick] table[x=t_s,y=y]{figdata/low-latency/16-protocols-ii-binary-exchange-protocols/gaps_events_A.csv};
\addplot[only marks,mark=|,mark size=4pt,omMeth,thick] table[x=t_s,y=y]{figdata/low-latency/16-protocols-ii-binary-exchange-protocols/gaps_events_B.csv};
\addplot[only marks,mark=*,mark size=2.5pt,omThm] table[x=t_s,y=y]{figdata/low-latency/16-protocols-ii-binary-exchange-protocols/gaps_events_arb.csv};
\draw[omThm,dashed] (axis cs:10.03,0.5) -- (axis cs:10.03,3.5);
\node[anchor=west,font=\scriptsize,omThm] at (axis cs:10.6,1.35) {both lines out: 17 messages, and 1 more};
\end{axis}
\end{tikzpicture}

\omcaption{Gaps on each line and on the arbitrated stream over 60 seconds of the simulator's feed (about 32\,000 packets), with
0.1\% independent loss per line, Gilbert--Elliott bursts, and outages of 60 milliseconds on each line overlapping by 25.
Each mark is one gap. Data: \texttt{fig\_gaps.py} (deterministic: the simulator is seeded).}\label{fig:ll:protocols-ii-binary-exchange-protocols:gaps}
\end{omfigure}

Arbitration costs little: the build's arbiter handles a packet in about \qty{1.3}{\nano\second}, a comparison of two sequence
numbers. What costs is the design around it: two network interfaces, two paths through two sets of switches, and a feed
handler that reads both without letting one line's delay hold up the other (chapter~18).

\begin{dated}{2026-09}{CME market data}\label{dat:ll:protocols-ii-binary-exchange-protocols:cme}
CME Group's MDP~3.0 market data (client systems documentation, consulted September 2026) is sent on Incremental Feeds A and B
over UDP, and CME ``strongly recommends that client systems process both''; any packet ``can arrive first on either feed''.
Its messages use Simple Binary Encoding, little-endian, ``optimized for low latency of encoding and decoding''; recovery uses
snapshot loops whose messages carry the last incremental sequence number processed.
\end{dated}

\section{Binary order entry and simple binary encoding}

\begin{definition}[Simple binary encoding]\label{def:ll:protocols-ii-binary-exchange-protocols:sbe}
\emph{Simple binary encoding}\index{simple binary encoding} (SBE) is the FIX Trading Community's standard for binary messages:
fields of native binary types at fixed positions described by an XML message schema, each message preceded by a small header
that gives the length of its fixed block, a template identifier, the schema's identifier and its version; the byte order is a
property of the schema, little-endian by default.
\end{definition}

Order entry is binary too at most venues where speed matters. The simulator's order-entry protocol (One Quant Book~10,
chapter~26) runs over SoupBinTCP, which Nasdaq's specification calls ``a lightweight point-to-point protocol, built on top of
TCP/IP sockets'' that delivers the server's sequenced messages in order ``even across underlying TCP/IP socket connection
failures'': its sequenced messages carry no explicit number, both sides count them, and a client that logs in again asks for
the next number it wants. It plays the part of the \omterm{def:ll:protocols-i-fix:session}{FIX session} of chapter~15 at a fraction of the cost: a new order is a
38-byte message with its fields at fixed offsets.

SBE is the FIX Trading Community's answer to the same needs, standardised in 2016: its specification states a ``preference
for fixed positions and fixed length fields, supporting direct access to data''. The build derives an SBE-style schema for the
simulator's feed messages from the same JSON schema, little-endian, with the eight-byte message header and natural field sizes
(the six-byte timestamp widened to eight), and generates a second set of flyweights from that XML alone.

\omcode{code/firm/wirecodec/sbe/feed_sbe.xml}{14}{23}{The Add Order message in the generated SBE-style schema.}\label{lst:ll:protocols-ii-binary-exchange-protocols:sbe}

\Cref{fig:ll:protocols-ii-binary-exchange-protocols:decode} compares the decoders on 60 seconds of line A. Book~1's decoder,
which copies each message and calls a \texttt{std::function}, takes about \qty{11}{\nano\second} per message; the generated
flyweights about \qty{5}{\nano\second}, whichever the byte order: a byte swap is one instruction and costs nothing
measurable here, while the copy and the indirect call cost half. Framing whole packets and dispatching every message costs
about \qty{7}{\nano\second} per message in C++ and \qty{8}{\nano\second} in Rust. At a million messages a second, decoding is
under 1\% of a core: the protocol was designed for that.

\begin{omfigure}
\begin{tikzpicture}
\begin{axis}[width=11.5cm,height=5.4cm,ybar,bar width=16pt,ymin=0,ymax=14,
  symbolic x coords={book1,flyweight,sbe,packets,rust},xtick=data,
  xticklabels={{Book~1\\ decoder},{flyweight,\\ big-endian},{SBE-style,\\ little-endian},{packets,\\ C++},{packets,\\ Rust}},
  xticklabel style={align=center,font=\footnotesize},enlarge x limits=0.12,ylabel={ns per message},
  tick label style={font=\small},label style={font=\small},grid=major,grid style={gray!25},
  nodes near coords,every node near coord/.append style={font=\scriptsize},point meta=explicit symbolic,
  table/col sep=comma]
\addplot[fill=omDef!60,draw=omDef] table[x=key,y=ns_per_msg,meta=ns_per_msg]{figdata/low-latency/16-protocols-ii-binary-exchange-protocols/measured_decode.csv};
\end{axis}
\end{tikzpicture}

\omcaption{Decoding the order messages of 60 seconds of the simulator's line A (about 32\,000 messages), median of 21 passes:
Book~1's decoder (a copy and a \texttt{std::function} call per message), the generated flyweights in both byte orders, and,
in the last two bars, every packet framed and every message dispatched. Measured on a laptop (Intel Core Ultra 7 155H) under WSL2,
no isolated cores; GCC~11 at \texttt{-O2}, Rust 1.97 release. Data: \texttt{bench\_wire.py}.}\label{fig:ll:protocols-ii-binary-exchange-protocols:decode}
\end{omfigure}

\section{Tutorial: generated codecs and two lossy lines}\label{tut:ll:protocols-ii-binary-exchange-protocols:tut}

\begin{tutorial}
\textbf{Goal.} Generate the codecs from the simulator's schema, hold them to Book~10's golden messages, record two lossy
lines, arbitrate them and measure. \textbf{End state:}
\cref{fig:ll:protocols-ii-binary-exchange-protocols:gaps,fig:ll:protocols-ii-binary-exchange-protocols:decode}, and green
tests in Python, C++ and Rust.
\begin{tutsteps}
\item \textbf{Generate} with \texttt{python gen\_wirecodec.py}: the C++ and Rust flyweights from
\path{code/firm/exchsim/schema.json}, the SBE-style XML schema, and the little-endian flyweights from that XML. The test runs
the generator again and requires no difference.
\item \textbf{Check against the golden messages.} The C++ test decodes every message of Book~10's golden fixtures (feed,
order entry in both directions, control), compares each with its row in the golden CSV files, re-encodes it with the writer
and compares the bytes, and sends every feed message through the SBE-style layout and back; the Rust test decodes the same
messages and the recorded MoldUDP64 packets.
\item \textbf{Record two lossy lines.} \texttt{ll\_lines.record(60, dir)} runs the simulator with the losses of
\cref{fig:ll:protocols-ii-binary-exchange-protocols:gaps} and writes both lines in the recorded-file format.
\item \textbf{Arbitrate} with \cref{lst:ll:protocols-ii-binary-exchange-protocols:arb}; \texttt{bench\_wire.py} checks the C++
counts against the Python reference, then measures the decoders.
\end{tutsteps}
\textbf{What to change next.} Remove the outage overlap and check that the arbitrated stream has no gap; put the two lines
behind one simulated switch (the same outage on both) and watch every gap pass through.
\end{tutorial}

\section{Build: the wire codecs}\label{bld:ll:protocols-ii-binary-exchange-protocols:wirecodec}

\begin{build}
\bfield{Purpose} Every byte the firm reads from or writes to the simulator: the feed handler (chapter~18) decodes the feed
with it, the order gateway (chapter~21) encodes and decodes order entry, and the capture tools (chapter~23) read recorded
files.
\bfield{Interface} Generated by \texttt{gen\_wirecodec.py} from \path{code/firm/exchsim/schema.json}: C++20
\texttt{firm::wire} with a reader \texttt{Feed<T>}, \texttt{In<T>}, \texttt{Out<T>}, \texttt{Ctl<T>} and a writer
\texttt{Writer} per message (\texttt{FeedAWriter} and so on), \texttt{dispatch\_feed}, \texttt{dispatch\_in}, \texttt{dispatch\_out},
\texttt{dispatch\_ctl}, \texttt{MoldHeader}, \texttt{for\_each\_block}, \texttt{csv} and \texttt{rewrite}; the SBE-style
\texttt{firm::sbe} readers and writers and \texttt{to\_sbe}; Rust \texttt{firm\_wirecodec} with the same readers,
\texttt{csv\_feed}, \texttt{csv\_in}, \texttt{csv\_out}, \texttt{csv\_ctl}, \texttt{MoldHeader} and \texttt{for\_each\_block}.
\bfield{Rules} No hand-written layout: every offset comes from the schema; generated files are committed and regenerated
without difference; readers copy nothing and allocate nothing; a message whose length does not match its type is refused by
the dispatcher.
\bfield{Acceptance tests} \texttt{code/firm/wirecodec/}: every message of Book~10's golden fixtures decoded to its golden
values and re-encoded to the same bytes in C++, decoded in Rust; every feed message round-tripped through the SBE-style layout;
the recorded MoldUDP64 fixture framed; the generator deterministic and its outputs current.
\bfield{Stretch} Generate the SoupBinTCP frames and a MoldUDP64 retransmission request; generate from the XML schema in Rust
as well; a codec for a real venue's published schema.
\end{build}

\begin{omsources}
\item Nasdaq, \emph{TotalView-ITCH 5.0}, \emph{MoldUDP64} and \emph{SoupBinTCP} specifications.
\item CME Group Client Systems Wiki, MDP~3.0: ``Incremental Feed Arbitration'', ``Simple Binary Encoding'', ``Market Data
Snapshot -- Full Recovery''.
\item FIX Trading Community, \emph{Simple Binary Encoding}, version 1.0 (GitHub).
\item E.~N.~Gilbert, ``Capacity of a burst-noise channel'', \emph{Bell System Technical Journal} 39(5), 1960; E.~O.~Elliott,
``Estimates of error rates for codes on burst-noise channels'', \emph{Bell System Technical Journal} 42(5), 1963.
\item The protocol description of the simulator, \path{code/firm/exchsim/PROTOCOL.md} (One Quant Book~10, chapter~26).
\end{omsources}

\section{Exercises}

\begin{exercise}[$\star$]\label{exo:ll:protocols-ii-binary-exchange-protocols:1}
A MoldUDP64 packet has sequence number 1\,000 and count 3; the next packet on the same line has sequence number 1\,007. Which
messages are missing, and what does a \omterm{def:ll:protocols-i-fix:session}{heartbeat} with sequence number 1\,010 then tell you?
\end{exercise}

\begin{exercise}[$\star$]\label{exo:ll:protocols-ii-binary-exchange-protocols:2}
Decode by hand the price of an Add Order whose bytes 32--35 are \texttt{00 0F 42 A4}. What is it in currency units?
\end{exercise}

\begin{exercise}[$\star$]\label{exo:ll:protocols-ii-binary-exchange-protocols:3}
Each line loses 0.2\% of packets independently. What fraction does the arbitrated stream lose? And if both lines pass through
one switch that drops 0.01\% of packets?
\end{exercise}

\begin{exercise}[$\star\star$]\label{exo:ll:protocols-ii-binary-exchange-protocols:4}
Compute the stationary loss rate of the simulator's Gilbert--Elliott parameters (good to bad 0.0005, bad to good 0.1, 80\%
lost in the bad state) and the mean length of a bad period in packets. Why does burstiness matter for retransmission?
\end{exercise}

\begin{exercise}[$\star\star$]\label{exo:ll:protocols-ii-binary-exchange-protocols:5}
Why is market data sent over UDP \omterm{def:ll:protocols-ii-binary-exchange-protocols:multicast}{multicast} rather than TCP, and what does a subscriber give up?
\end{exercise}

\begin{exercise}[$\star\star$]\label{exo:ll:protocols-ii-binary-exchange-protocols:6}
The SBE-style layout widens the six-byte timestamp to eight bytes. What does that cost in bytes for an Add Order, and what does
it save?
\end{exercise}

\begin{exercise}[$\star\star\star$]\label{exo:ll:protocols-ii-binary-exchange-protocols:7}
\emph{Coding.} Extend \texttt{gen\_wirecodec.py} to generate the SoupBinTCP frames of the schema's \texttt{soup\_client} and
\texttt{soup\_server} protocols, and test them on Book~10's \texttt{fixture\_soup.bin}.
\end{exercise}

\begin{exercise}[$\star\star\star$]\label{exo:ll:protocols-ii-binary-exchange-protocols:8}
\emph{Find the flaw.} ``Our feed handler reads line A; if it sees a gap it switches to line B until the next gap.''
\end{exercise}

\section{Problem: Two Copies of Everything}

\begin{problem}[{Weekend problem --- how often redundancy fails, and what recovery costs}]\label{pb:ll:protocols-ii-binary-exchange-protocols:1}
A feed publishes 50\,000 packets a second over a 6.5-hour session on two lines. Assume each line loses a packet with probability
$10^{-4}$, independently; retransmission serves gaps of up to 1\,000 messages in a round trip of \qty{0.5}{\milli\second}; a
snapshot of the instrument arrives every second and takes \qty{20}{\milli\second} to receive; applying a buffered message costs
\qty{8}{\nano\second} (a little above the measured cost of framing and dispatching one).

\medskip\noindent\textbf{Part I --- Independent loss.}
\begin{enumerate}
\item How many packets does a session carry, and how many does one line lose?
\item What is the probability that the arbitrated stream loses a packet, and how many losses is that per session?
\item In the simulation of \cref{fig:ll:protocols-ii-binary-exchange-protocols:gaps}, why are there any arbitrated gaps at all?
\item What does each arbitrated loss cost with retransmission?
\end{enumerate}

\medskip\noindent\textbf{Part II --- Correlation.}
\begin{enumerate}[resume]
\item A shared switch drops $10^{-6}$ of packets on both lines. What is the arbitrated loss rate now (\cref{prop:ll:protocols-ii-binary-exchange-protocols:both}), and how many losses per session?
\item What stationary loss rate do the simulator's burst parameters give each line, and what arbitrated rate if the lines are
independent?
\item Why do bursts make retransmission less useful even when the average loss rate is unchanged?
\item Name three common causes a firm can remove, and one it cannot.
\end{enumerate}

\medskip\noindent\textbf{Part III --- Recovery.}
\begin{enumerate}[resume]
\item When is retransmission not enough?
\item How long does recovery from the \omterm{def:ll:protocols-ii-binary-exchange-protocols:channels}{snapshot channel} take on average, and at worst?
\item How many incremental messages arrive meanwhile, and how long does applying them take?
\item Which buffered messages must be discarded after the snapshot, and how does the handler know?
\end{enumerate}

\medskip\noindent\textbf{Part IV --- The verdict.}
\begin{enumerate}[resume]
\item State the \emph{named result}: the arbitrated loss probability with independent and with correlated losses, and the
recovery time from the \omterm{def:ll:protocols-ii-binary-exchange-protocols:channels}{snapshot channel}.
\item What should a strategy do with its view of the book during a recovery?
\item Would a snapshot every 100 milliseconds help? What does it cost the venue and the subscriber?
\item Why must the arbiter never wait for the slower line?
\item How would you test the recovery path before production?
\item What should be logged for every gap?
\item What does the measured cost of decoding say about where the latency of a feed handler lies?
\item In one sentence: what does a second line buy?
\end{enumerate}
\end{problem}

\section{Interview questions}

\begin{interviewq}[$\star$ \iqroles{developer}]\label{iq:ll:protocols-ii-binary-exchange-protocols:1}
Why do exchanges publish market data over UDP \omterm{def:ll:protocols-ii-binary-exchange-protocols:multicast}{multicast}, and how does a receiver detect a lost packet?
\end{interviewq}

\begin{interviewq}[$\star\star$ \iqroles{developer}]\label{iq:ll:protocols-ii-binary-exchange-protocols:2}
Implement A/B \omterm{def:ll:protocols-ii-binary-exchange-protocols:arbitration}{line arbitration}. What state do you keep, and what do you do on a gap?
\end{interviewq}

\begin{interviewq}[$\star\star$ \iqroles{developer}]\label{iq:ll:protocols-ii-binary-exchange-protocols:3}
What is a flyweight decoder, and why is it faster than decoding into structs?
\end{interviewq}

\begin{interviewq}[$\star\star$ \iqroles{developer}]\label{iq:ll:protocols-ii-binary-exchange-protocols:4}
You joined the feed in the middle of the day. How do you build a correct book?
\end{interviewq}

\begin{interviewq}[$\star\star$ \iqroles{developer}]\label{iq:ll:protocols-ii-binary-exchange-protocols:5}
Big-endian or little-endian on the wire: does it matter for latency? What does?
\end{interviewq}

\begin{interviewq}[$\star\star\star$ \iqroles{developer, researcher}]\label{iq:ll:protocols-ii-binary-exchange-protocols:6}
Your two lines are ``redundant'' but you see simultaneous gaps on both several times a day. What do you investigate?
\end{interviewq}
